%%%%%%%%%%%%%%%%%%%%%%%%
\documentclass[epjST]{svjour}

\voffset -0.7cm
\hoffset -0.7cm
\usepackage{amsmath,amssymb,bbm}
\usepackage{float}
\usepackage{imakeidx}
\usepackage[colorlinks=true,linkcolor=blue,citecolor=blue,urlcolor=magenta]{hyperref}
\usepackage{footnote,comment}
\usepackage{doi}
\usepackage{enumitem}
\usepackage{slashed}
\usepackage{xcolor,xstring,xparse}
\usepackage{graphicx}
\usepackage{rotating}
\usepackage{pdflscape}
\usepackage{aastex_hack}
\usepackage{appendix}
\usepackage{booktabs}

\makesavenoteenv{tabular}
\makesavenoteenv{table}
\makesavenoteenv{figure}

% Make Orcid icon
\newcommand{\orcidicon}{\includegraphics[width=0.32cm]{orcid.pdf}}
\newcommand{\orc}[1]{\href{https://orcid.org/#1}{\orcidicon}}

% Author Orcid ID: Define per author
\newcommand{\orcJR}{0000-0001-8217-1484}
\newcommand{\orcCTY}{0000-0001-5038-8427}
\newcommand{\orcAJS}{0000-0001-5474-2649}
\newcommand{\orcMF}{0000-0003-2704-6474}
\newcommand{\orcJB}{0000-0002-2289-4856}
\newcommand{\orcSE}{0000-0001-5884-5047}
\newcommand{\orcCG}{0000-0001-9985-1822}
\newcommand{\orcWP}{0000-0003-0793-3041}
\newcommand{\orcLL}{0000-0002-1478-5254}

% Useful macros for equations and units
\newcommand{\ts}{\textstyle}
\newcommand*{\TeV}{\text{\,TeV}}
\newcommand*{\GeV}{\text{\,GeV}}
\newcommand*{\MeV}{\text{\,MeV}}
\newcommand*{\keV}{\text{\,keV}}
\newcommand*{\eV}{\text{\,eV}}
\newcommand*{\meV}{\text{\,meV}}
\newcommand*{\Msun}{\mathrm{M}_{\odot}}
\newcommand*{\bb}{\boldsymbol}
\newcommand*{\beqn}{\begin{equation}}
\newcommand*{\eeqn}{\end{equation}}
\newcommand{\beql}[1]{\begin{equation} \label{#1}}
\newcommand{\eeql}[1]{\label{#1} \end{equation} }  
\newcommand{\req}[1]{Eq.\,(\ref{#1})}
\newcommand{\rf}[1]{Fig.\,{\ref{#1}}}
\newcommand{\rt}[1]{Table~{\ref{#1}}}
\newcommand{\rsec}[1]{Sec.\,{\ref{#1}}}
\newcommand{\rchap}[1]{Ch.\,{\ref{#1}}}
\newcommand{\rapp}[1]{Appendix~{\ref{#1}}}
\newcommand{\mydoi}[2]{\href{http://dx.doi.org/#2}{#1}}
\newcommand{\E}{\mathrm{e}}
\newcommand{\ie}{{\em i.e.\/}}  
\newcommand{\eg}{{\em e.g.\/}}  
\newcommand{\ms}[1]{\rule[-#1mm]{0mm}{0mm}}
\newcommand{\grad}{\operatorname{grad}}
\newcommand{\diag}{\mathrm{diag}}

% Andrew's commands
\newcommand{\cccite}[1]{Published in Ref.~\cite{#1} under the \href{https://creativecommons.org/licenses/by/4.0/}{CC BY 4.0} license}
\newcommand{\radapt}[1]{Adapted from Ref.~\cite{#1}}
\newcommand{\para}[1]{\paragraph{#1}\hfill\break\noindent}
\NewDocumentCommand{\allcite}{m o}{\nocite{#1}\hyperlink{cite.#1}{\StrBefore{#1}{:}\IfNoValueTF{#2}{ et. al. }{ and #2 }(\StrBehind{#1}{:}[\temp]\StrLeft{\temp}{4})}}
\NewDocumentCommand{\aucite}{m o}{\nocite{#1}\hyperlink{cite.#1}{\StrBefore{#1}{:}\IfNoValueTF{#2}{ }{ and #2 }(\StrBehind{#1}{:}[\temp]\StrLeft{\temp}{4})}}
\NewDocumentCommand{\allcitep}{m o}{\nocite{#1}\hyperlink{cite.#1}{\StrBefore{#1}{:}\IfNoValueTF{#2}{ et. al. }{ and #2 }(preprint \StrBehind{#1}{:}[\temp]\StrLeft{\temp}{4})}}

% Useful macros for annotation
\newcommand*{\xred}{\color{red}}
\newcommand*{\xblue}{\color{blue}}
\newcommand*{\xgreen}{\color{green}}

% Struts for tables 
\newcommand\Tstrut{\rule{0pt}{2.6ex}}
\newcommand\Bstrut{\rule[-0.9ex]{0pt}{0pt}}
\newcommand{\TBstrut}{\Tstrut\Bstrut}

% Equations numbered by section
%\numberwithin{equation}{section}

\begin{document}

\title{Short note on spin magnetization in quark-gluon plasma
    }

\author{
    Andrew~Steinmetz${}^1$\orc{\orcAJS},
    Johann~Rafelski${}^1$\orc{\orcJR}
    }

\institute{
    ${}^1$Department of Physics, The University of Arizona, Tucson, AZ, 85721, USA\\
    }

\abstract{
    We outline the theory of spin magnetization of the primordial quark-gluon plasma (QGP) that existed shortly after the Big Bang, focusing on a magnetized fermion gas of light quarks and electrons within temperatures of \(150\) MeV to \(500\) MeV. Using a grand partition function approach and evaluating magnetized Fermi-Dirac integrals, we calculate the magnetization and estimate that a fully polarized up quark gas could generate cosmic magnetic fields of approximately \(10^{16}\) Tesla. This field strength is significantly larger than the critical Schwinger field and exceeds magnetar surface fields by over three orders of magnitude. This suggest that even a weakly polarized quark gas would have a profound impact on the early Universe.
    }

\maketitle
%%%%%%%%%%%%%%%%%%%%%%%%

%%%%%%%%%%%%%%%%%%%%%%%%
\section{Introduction}
\label{sec:introduction}
%%%%%%%%%%%%%%%%%%%%%%%%
Quark-gluon plasma (QGP) is a state of matter that existed microseconds after the Big Bang, where quarks and gluons were not confined within hadrons but existed as a free deconfined plasma~\cite{Letessier:2002ony,Rafelski:2015cxa,Rafelski:2023emw,Rafelski:2024fej,Shuryak:2004cy,HotQCD:2014kol}. QGP is also produced in high-energy heavy-ion collisions at the Relativistic Heavy Ion Collider (RHIC) and Large Hadron Collider (LHC)~\cite{Rafelski:1982pu,STAR:2005gfr,Palni:2024wdy,Grayson:2024okq,Mu:2025gtr}. While heavy-ion collisions provide our best experimental probe of the properties of primordial QGP, there are differences between the QGP produced in the early Universe (nearly baryon-symmetric) versus QGP produced in laboratory heavy-ion collisions (baryon-rich) as well as differing lepton abundance due to cosmic charge neutrality~\cite{Yang:2024ret}. Primordial QGP also differs from dense quark matter in astrophysical compact objects where the quark chemical potential far exceeds the electron chemical potential~\cite{Ghosh:2025sjn}.

Understanding the properties of QGP and cosmic magnetism is crucial for determining the conditions of the primordial Universe and the formation of matter as we observe it today~\cite{Uzan:2010pm,Uzan:2024ded,Grasso:2000wj,Widrow:2002ud,Kandus:2010nw,Ghosh:2024owm}.
\begin{table}[h]
    \centering
    \caption{Properties of select particles}
    \label{tab:particle_properties}
    \begin{tabular}{@{}lllll@{}}
        \toprule
        \textbf{Particle} & \textbf{Mass} \([\approx\mathrm{MeV}]\) & \textbf{Charge} & \textbf{Magneton} \([\mu/\mu_{B}]\) & \(\mathfrak{g}\) \textbf{dof.} \\ 
        \midrule
        Electron    \((e)\) & 0.511 & \(-1\)    & \(-1\)         & 2 \\
        Up          \((u)\) & 2.2   & \(+2/3\)  & \(+0.155\)     & 6 \\
        Down        \((d)\) & 4.7   & \(-1/3\)  & \(-0.0362\)    & 6 \\
        Strange     \((s)\) & 96    & \(-1/3\)  & \(-0.00177\)   & 6 \\
        Charm       \((c)\) & 1270  & \(+2/3\)  & \(+0.000268\)  & 6 \\ 
        \bottomrule
    \end{tabular}
\end{table}
In Table~\ref{tab:particle_properties}, we list the properties of select particles relevant to the QGP era. The magnetic moment \(\mu\) is given in units of the Bohr magneton \(\mu_{B}\equiv e\hbar/2m_{e}\approx0.5927\ \mathrm{MeV}^{-1}\). The degrees of freedom (dof.) \(\mathfrak{g}=n_{S}\times n_{C}\) is the number of spin \(n_\mathrm{S}\) and color \(n_\mathrm{C}\) states available to the particle. We evaluated the magneton with gyromagnetic factor \(g=2\).

The electron-positron and light-quark gases were magnetically the most relevant particles in the QGP era due to their charge and low mass. Their number densities \(n_{i}\) would approximate a massless gas following the expression
\begin{align}
    n_{i}(T) = \frac{3\mathfrak{g}_{i}}{4\pi^{2}}\zeta(3)T^{3}\,,
\end{align}
where \(\zeta(3)\approx1.202\) is the Zeta function. If the dense gas of up (anti)quarks and electron(positrons) was fully polarized at a temperature of \(300\) MeV prior to hadronization, this would correspond to an estimated cosmic magnetization flux strength \(M = \mathcal{M}/V\) of
\begin{align}
    M_{u\bar{u}} = \mathcal{M}_{u\bar{u}}/V = \mu_{u}(n_{u}+n_{\bar{u}}) = \frac{\mu_{u}}{\mu_{B}}\mu_{B}\frac{3\mathfrak{g}_{u}}{2\pi^{2}}\zeta(3)T^{3} & \sim 1.391\times10^{16}\ \mathrm{T}\,,\\
    M_{e\bar{e}} = \frac{\mu_{e}}{\mu_{u}} \frac{\mathfrak{g}_{e}}{\mathfrak{g}_{u}} M_{u\bar{u}} = & \sim 2.991\times10^{16}\ \mathrm{T}\,,
\end{align}
where \(\mathcal{M}\) is the magnetization. While this is a overly simplified calculation, it illustrates the scale at which spin magnetization may be relevant. These values can also be understood as \(\sim 10^{6}\) times the critical magnetic field strength \(\mathcal{B}_\mathrm{C}\equiv m_{e}^{2}/e\ \approx 0.862\ \mathrm{MeV}^{2}\) and \(\sim 10^5\) times stronger than the upper estimated surface field strength of magnetars~\cite{Kaspi:2017fwg}. This suggests that even a weakly polarized primordial gas would have a significant impact on the early Universe, with light quarks contributing on par with leptons. We can understand this as despite having a smaller magneton, the light quarks have a larger number density, due to added color degrees of freedom, and therefore contribute more to spin magnetization than otherwise expected.

In Ref.~\cite{Steinmetz:2023nsc,Steinmetz:2023ucp}, we explored the magnetic properties of the magnetized electron-positron plasma around the period of Big Bang Nucleosynthesis~\cite{Grayson:2023flr,Grayson:2024uwg} as well as neutrino flavor~\cite{Rafelski:2023zgp}. We expand those efforts to now consider the thermodynamic properties of a magnetized QGP in the primordial Universe, focusing on the interplay between quarks, leptons, and magnetic fields. We provide an theoretical outline and point to where future efforts may be directed. The presence of strong magnetic fields in the primordial Universe could have significantly affected the equilibrium properties of Standard Model particles in the earliest moments after the Big Bang~\cite{Durrer:2013pga,Subramanian:2015lua}. Such magnetic fields have long been thought to be connected to baryon asymmetry~\cite{Vachaspati:1991nm,Baym:1995fk}. Chiral magnetism in QGP has also been studied~\cite{Fukushima:2008xe,Boyarsky:2011uy,Bali:2011qj}.

%%%%%%%%%%%%%%%%%%%%%%%%
\section{Magnetization of a polarized gas}
\label{sec:magnetization}
%%%%%%%%%%%%%%%%%%%%%%%%
We consider a free but magnetized fermion gas in the temperature range \(500\MeV>T>150\MeV\) composed of light-quark species \(q \in {u,d}\) and electrons. As the magneton scales with \(\mu \propto 1/m\), these species are magnetically the most relevant due to their lighter masses (see Table~\ref{tab:particle_properties}) and consequently larger magnetic moments. For relativistic species~\cite{Elze:1980er}, under conditions of thermal and chemical equilibrium (as was the case in the primordial Universe) the chemical potential \(\eta\) of each particle is opposite in sign to that of its antiparticle
\begin{align}
    \eta_{q}=T\ln\lambda_{q}\,,\qquad
    \lambda_{q}=1/\lambda_{\bar{q}}\,,\qquad
    \eta_{q}=-\eta_{\bar{q}}\,,
\end{align}
where \(\lambda\) is the fugacity. The magnetic dipole of a particle is also opposite in sign to its antiparticle $\mu_{i}=-\mu_{\bar{i}}$ as charge is flipped. Any deviation from this condition would represent a violation of CPT~\cite{Colladay:1996iz,Bluhm:1997ci,BASE:2016yuo}. During this period, particle-antiparticle pairs of (anti)quarks were freely produced and annihilated through photon-mediated processes, represented by $q+\bar{q}\rightleftharpoons2\gamma$. Since the Universe expanded very slowly compared to collision reaction times during the QGP epoch~\cite{Rafelski:2023emw,Yang:2024ret}, the expansion can be considered adiabatic, conserving entropy.

Accounting for the internal energy $U$ of magnetized QGP, including the energies of neutrinos~\cite{Birrell:2014ona}, involves the following properties: 
\begin{itemize}
    \item[(a)] The energy of adding or removing a baryon $\eta_{B}B$,
    \item[(b)] the energy of adding or removing a lepton $\eta_{\ell}(N_{\ell}-N_{\ell})$ with $\ell\in {e,\nu}$, 
    \item[(c)] the magnetic energy $\mathcal{M}\mathcal{B}$ where $\mathcal{M}$ is the net magnetization and $\mathcal{B}$ is the magnetic field strength and
    \item[(d)] the electromagnetic energy density generated by the external magnetic field.
\end{itemize}
Magnetic field strength $\mathcal{B}$ should not be confused with baryon number $B$. Additionally, the internal energy and chemical potential of QGP is sensitive to electric fields which cause a statistical Coulomb distortion~\cite{Sigl:1996dm,Letessier:2002ony}. For the purpose of clarity and to maintain focus on magnetism, electrical fields will be omitted. See Ref.~\cite{Grayson:2022asf,Shovkovy:2022bnd,Ghosh:2024fkg} for discussion on QGP electrical conductivity and EM response.

The dependency of $U$ on $\mathcal{M}$ reflects that $\mathcal{B}$ is the incremental energy cost to change the magnetization by flipping the spin of a particle~\cite{Bali:2014kia}. Therefore, this makes magnetization $\mathcal{M}$ an extensive property of the system which changes which particle number. We see this explicitly by writing the magnetization as the sum over all particles $i\in{1,\ldots,k}$
\begin{align}
    \label{eq:dipole}
    \mathcal{M} = \sum_{i=1}^{k}(\mu_{i}N_{i}^{\uparrow} + \mu_{\bar{i}}N_{\bar{i}}^{\uparrow} - \mu_{i}N_{i}^{\downarrow} - \mu_{\bar{i}}N_{\bar{i}}^{\downarrow})\,,\qquad
    N_{i} = N_{i}^{\uparrow} + N_{i}^{\downarrow}\,,
\end{align}
where $\mu_{i}$ is the magnetic dipole moment per particle. The $\uparrow\downarrow$ notation refers to spin-up $(\uparrow)$ and spin-down $(\downarrow)$ states along the direction of the external field. Therefore, $N_{i}^{\uparrow\downarrow}$ refers to the $i$-th constituent population number in either spin-up or spin-down orientation. The signs of each term in \req{eq:dipole} arises from the sign of the spin eigenvalue. While \req{eq:dipole} presumably includes contributions from each particle with a magnetic dipole, we expect the magnetization to be dominated by electrons(positrons) and the lightest quarks due to their charge and low mass therefore we sum over $i\in{u,d,e}$
\begin{align}
    \notag\mathcal{M} \approx &+|\mu_{u}|(N_{u}^{\uparrow}-N_{\bar{u}}^{\uparrow})-|\mu_{u}|(N_{u}^{\downarrow}-N_{\bar{u}}^{\downarrow})\\
    \notag &-|\mu_{d}|(N_{d}^{\uparrow}-N_{\bar{d}}^{\uparrow})+|\mu_{d}|(N_{d}^{\downarrow}-N_{\bar{d}}^{\downarrow})\\
    \label{eq:dipole2}
    &-|\mu_{e}|(N_{e}^{\uparrow}-N_{\bar{e}}^{\uparrow})+|\mu_{e}|(N_{e}^{\downarrow}-N_{\bar{e}}^{\downarrow})\,.
\end{align}

We recognize that \req{eq:dipole2} contains terms representing asymmetry in the spin alignment though we can organize them in two different ways: (a) We group terms of the same spin alignment or (b) we group terms of matter and antimatter. The second approach may allow definition of spin-asymmetry in terms of conserved quantities characterizing spin angular momentum. We define net spin-asymmetry numbers $\delta_{i}^{\uparrow\downarrow}$ and write
\begin{align}
    \delta_{i}^{\uparrow\downarrow} &\equiv N_{i}^{\uparrow\downarrow}-N_{\bar{i}}^{\uparrow\downarrow}\,,\\
    \mathcal{M} &= 
    +|\mu_{u}|(\delta_{u}^{\uparrow}-\delta_{u}^{\downarrow})
    -|\mu_{d}|(\delta_{d}^{\uparrow}-\delta_{d}^{\downarrow})
    -|\mu_{e}|(\delta_{e}^{\uparrow}-\delta_{e}^{\downarrow})\,.
\end{align}
The net spin-asymmetry warrants is the asymmetry of particles and antiparticles of the same spin. Therefore $\delta_{u}^{\uparrow}\neq0$ represents a situation where there are more up quarks than up antiquarks in the spin-up $\uparrow$ state.

%%%%%%%%%%%%%%%%%%%%%%%%
\section{Magnetized grand partition function}
\label{sec:partition}
%%%%%%%%%%%%%%%%%%%%%%%%
The partition function allows us to calculate various thermodynamic quantities by taking appropriate derivatives of $\mathcal{F}$. In the temperature range considered $(500\MeV>T>150\MeV)$, the lightest quarks act as essentially massless particles with only the strange quark requiring significant mass corrections. It is worth remarking on the uniqueness of the situation: As magnetic moment scales inverse with mass, it is the particles which are most massless in character which contribute most to magnetization.

The relevant contributions to the magnetized primordial plasma arise from the quarks, gluons, leptons, and the vacuum. The grand potential in terms of the grand partition function $\ln\mathcal{Z}$ is
\begin{align}
    \label{eq:parts}
    \mathcal{F} &= -T\ln\mathcal{Z}\,,\\
    \ln\mathcal{Z}_{\mathrm{total}} &=
    \ln\mathcal{Z}_{\mathrm{quarks}} +
    \ln\mathcal{Z}_{\mathrm{gluons}} +
    \ln\mathcal{Z}_{\mathrm{vac.}} + 
    \ln\mathcal{Z}_{\mathrm{leptons}}+\ldots 
\end{align}

We consider a homogeneous magnetic field domain defined along the $z$-axis with magnetic field magnitude $\mathcal{B}$. The volume $V=L^{3}$ is not necessarily infinite and is to be considered the size of the homogeneous domain such that $\partial\mathcal{B}_{i}/\partial x_{j}\approx0$ for \(i,j \in {1,2,3}\). For a fermion species of charge $Q$, mass $m$, and g-factor $g$, the energy eigenvalues of the magnetized particles is given by~\cite{Steinmetz:2018ryf}
\begin{align}
    \label{eq:energystates}
    E(p_{z},n,s)=\sqrt{m^{2}+p_{z}^{2}+2|Q|\mathcal{B}\left(n+\frac{1}{2}-\frac{g}{2}s\right)}\,,
\end{align}
where $E$ are the relativistic Landau energy eigenvalues. The micro-state energies depend on longitudinal momentum \(p_{z}\), spin $s\in\pm1/2$, and orbital $n\in0,1,2,3,\ldots$ Landau quantum number. It is helpful to introduce a spin-dependent auxiliary mass $m_{T}(s)$ via
\begin{align}
    \label{eq:spinmass}
    m^{2}(s) &\equiv m^{2} - |Q|\mathcal{B}\,g s\,.
\end{align}

The power and utility of the partition function in statistical systems is found by examining the Fermi integral in various limits and expansions. We define the Fermi-Dirac distribution noting that fugacity \(\lambda\) is related to chemical potential in the usual way via \(\eta = T\ln\lambda\). Thus,
\begin{align}
    F\left(E - \sigma\eta\right) = \frac{1}{e^{(E - \sigma\eta)/T} + 1}\,.
\end{align}
We can express \req{eq:parts} by utilizing Euler-Maclaurin integration (see details in Ref.~\cite{Steinmetz:2023nsc,Rafelski:2024fej}) writing the partition function in spherical coordinates \(d\mathbf{p}^{3}=4\pi p^{2}dp\). We substitute coordinates and integrate by parts yielding
\begin{align}
    \label{eq:partition_byparts}
    \ln\mathcal{Z} &= \frac{2 n_\mathrm{C}V}{(2\pi)^{2}} \sum_{s}^{\pm1/2}\sum_{\sigma}^{\pm1}\int_{0}^{\infty} \frac{dp}{3T} \, \frac{p^4}{E}F\left(E - \sigma\eta\right)\,.
\end{align}
The form of the partition function expressed by \req{eq:partition_byparts} more directly lets us evaluate thermodynamic quantities in terms of Fermi integrals~\cite{Elze:1980er,Birrell:2024bdb}. \req{eq:partition_byparts} also sums over (anti)matter states \(\sigma \in {\pm1}\). However, integrating over momentum is not an ideal description as relativistic expansions in momentum yield series that are only semi-convergent. These infrared divergences in thermal field theory are associated with long-wavelength (low-momentum) modes that are not properly accounted for in the perturbative expansion~\cite{Gross:1980br,Linde:1980ts,Bellac:2011kqa,Kapusta:2006pm}.

%%%%%%%%%%%%%%%%%%%%%%%%
\subsection{Dimensionless change of variables}
\label{sec:dimensionless}
%%%%%%%%%%%%%%%%%%%%%%%%
To simplify the integration process, we introduce dimensionless variables by normalizing relevant physical quantities with the temperature \( T \). This approach renders the equations dimensionless and highlights the thermal contributions explicitly. The dimensionless variables are defined as
\begin{align}
    \label{eq:dimensionless_variables}
    p_{T} = \frac{p}{T}, \qquad E_{T}(p_{T},s) = \frac{E(p,s)}{T}, \qquad \eta_{T} = \frac{\eta}{T}\,, \qquad m_{T}(s) = \frac{m(s)}{T}\,.
\end{align}
This defines momentum-like \(p_{T}\), energy-like \(E_{T}\), chemical potential-like \(\eta_{T}\) and mass-like \(m_{T}\) parameters. Using the relativistic dispersion relation, the dimensionless energy \( E_{T} \) can be expressed in terms of \( p_{T} \) and \( m_{T} \)
\begin{align}
    E_{T} = \frac{E}{T} = \sqrt{p_{T}^{2} + m_{T}^{2}}\,.
\end{align}
The differential \( dp_{T} \) and \( dE_{T} \) transform as
\begin{align}
    dp = T \, dp_{T}\,,\qquad p_{T}dp_{T} = E_{T} dE_{T}\,,
\end{align}
and the limits of integration change accordingly
\begin{equation}
    p_{T} = 0 \quad \Rightarrow \quad E_{T} = m_{T}, \quad p_{T} \to \infty \quad \Rightarrow \quad E_{T} \to \infty.
\end{equation}
Substituting these dimensionless variables and differentials into the partition function \( \ln\mathcal{Z} \), we obtain expressions for both momentum-like \(p_{T}\) integration and energy-like \(E_{T}\) integration
\begin{align}
    \label{eq:dimensionless_partition}
    \ln\mathcal{Z} 
    &= \frac{2 n_\mathrm{C} V}{(2\pi)^{2}} \frac{T^{3}}{3} \sum_{s}^{\pm1/2}\sum_{\sigma}^{\pm1} \int_{0}^{\infty} dp_{T} \, \frac{p_{T}^{4}}{\sqrt{p_{T}^{2} + m_{T}^{2}}} \, F\left(\sqrt{p_{T}^{2} + m_{T}^{2}} - \sigma\eta_{T}\right)\,,\\
    \label{eq:dimensionless_partition2}
    &= \frac{2 n_\mathrm{C} V}{(2\pi)^{2}} \frac{T^{3}}{3} \sum_{s}^{\pm1/2}\sum_{\sigma}^{\pm1} \int_{m_{T}}^{\infty} dE_{T} \, (E_{T}^2-m_{T}^2)^{3/2} F\left(E_{T} - \sigma\eta_{T}\right)\,.
\end{align}
In this formulation, it is evident that the logarithm of the partition function scales as \( \ln\mathcal{Z} \propto T^{3} \), consistent with the expected thermodynamic behavior for a relativistic gas in three spatial dimensions.

%%%%%%%%%%%%%%%%%%%%%%%%
\subsection{Evaluation of magnetization from the dimensionless partition function}
\label{sec:magnetization_evaluation}
%%%%%%%%%%%%%%%%%%%%%%%%
Given the dimensionless form of the partition function \req{eq:dimensionless_partition}, we proceed to evaluate the magnetization. We emphasize that the dimensionless mass \(m_{T}(\mathcal{B},s)\) depends on the magnetic field and spin via \req{eq:spinmass}. Taking the derivative of the free energy \(\mathcal{F} = -T \ln\mathcal{Z}\) with respect to the magnetic field \(\mathcal{B}\), we obtain the magnetization
\begin{align}
    \label{eq:magnetization_def}
    \mathcal{M} = \left( \frac{\partial \mathcal{F}}{\partial \mathcal{B}} \right) = -T \left( \frac{\partial \ln\mathcal{Z}}{\partial \mathcal{B}} \right)\,.
\end{align} 
Since \(\ln\mathcal{Z}\) depends on \(\mathcal{B}\) solely through \(m_{T}\), we apply the chain rule
\begin{align}
    \label{eq:chain_rule}
    \frac{\partial \ln\mathcal{Z}}{\partial \mathcal{B}} = \frac{\partial \ln\mathcal{Z}}{\partial m_{T}} \frac{\partial m_{T}}{\partial \mathcal{B}}\,,\qquad
    \frac{\partial m_{T}}{\partial \mathcal{B}} = -\frac{g|Q|s}{2m_{T}T^2}\,.
\end{align}
Taking the derivative of the partition function with respect to \(m_{T}\), we write
\begin{equation}
    \frac{\partial \ln \mathcal{Z}}{\partial m_{T}} = \frac{2 n_\mathrm{C} V T^{3}}{3 (2\pi)^{2}} \sum_{s}^{\pm1/2}\sum_{\sigma}^{\pm1} \int_{0}^{\infty} dp_{T} \, \frac{\partial}{\partial m_{T}} \left( \frac{p_{T}^{4}}{\sqrt{p_{T}^{2} + m_{T}^{2}}} F\left(\sqrt{p_{T}^{2} + m_{T}^{2}} - \sigma \eta_{T} \right) \right).
\end{equation}
Given that \(E_{T} = \sqrt{p_{T}^{2} + m_{T}^{2}}\), then, \(\partial E_{T} / \partial m_{T} = m_{T} / E_{T}\). The derivative of the integrand is computed using the product and chain rules
\begin{equation}
    \frac{\partial}{\partial m_{T}} \left( \frac{p_{T}^{4}}{E_{T}} F(E_{T} - \sigma\eta_{T}) \right) = \frac{p_{T}^{4}}{E_{T}} \frac{\partial F}{\partial E_{T}} \frac{\partial E_{T}}{\partial m_{T}} + F(E_{T} - \sigma\eta_{T}) \frac{\partial}{\partial m_{T}} \left( \frac{p_{T}^{4}}{E_{T}} \right).
\end{equation}
Hereafter we write \(F'=\partial F/\partial E_{T}\). The second term evaluates to
\begin{equation}
    \frac{\partial}{\partial m_{T}} \left( \frac{p_{T}^{4}}{E_{T}} \right) = -\frac{p_{T}^{4} m_{T}}{E_{T}^{3}}.
\end{equation}
Substituting these results back, the integrand becomes
\begin{equation}
    \frac{\partial}{\partial m_{T}} \left( \frac{p_{T}^{4}}{E_{T}} F(E_{T} - \sigma\eta_{T}) \right) = \frac{p_{T}^{4} m_{T}}{E_{T}^{2}} F'(E_{T} - \sigma\eta_{T}) - \frac{p_{T}^{4} m_{T}}{E_{T}^{3}} F(E_{T} - \sigma\eta_{T}).
\end{equation}
Replacing \(E_{T}\) with \(\sqrt{p_{T}^{2} + m_{T}^{2}}\), the derivative of \(\ln \mathcal{Z}\) is
\begin{align}
    \frac{\partial \ln \mathcal{Z}}{\partial m_{T}} &= \frac{2 n_\mathrm{C} V T^{3}}{3 (2\pi)^{2}} \sum_{s}^{\pm1/2}\sum_{\sigma}^{\pm1} \int_{0}^{\infty} dp_{T} \, p_{T}^{4} m_{T}
    \left( \frac{F'\left( \sqrt{p_{T}^{2} + m_{T}^{2}} - \sigma\eta_{T} \right)}{p_{T}^{2} + m_{T}^{2}} - \frac{F\left( \sqrt{p_{T}^{2} + m_{T}^{2}} - \sigma\eta_{T} \right)}{(p_{T}^{2} + m_{T}^{2})^{3/2}} \right).
\end{align}
This result provides the explicit form of \(\partial \ln \mathcal{Z}/\partial m_{T}\) in terms of \(F\) and its derivative \(F'\), with all dependencies on \(m_{T}\) and \(p_{T}\) made explicit.

Given that \( F(x) = 1/(e^{x} + 1) \) is the Fermi-Dirac distribution, its derivative is
\begin{equation}
F'(x) = \frac{dF}{dx} = -\frac{e^{x}}{(e^{x} + 1)^2} = -F(x)\left[1 - F(x)\right].
\end{equation}
We replace \( F'(x) \) in the expression for the derivative of the integrand yielding
\begin{align}
    \notag
    &\int_{0}^{\infty} dp_{T} \, p_{T}^{4} m_{T} \left( -\frac{F(E_{T} - \sigma\eta_{T})\left[1 - F(E_{T} - \sigma\eta_{T})\right]}{E_{T}^{2}} - \frac{F(E_{T} - \sigma\eta_{T})}{E_{T}^{3}} \right) =\\*
    &\int_{m_{T}}^{\infty} dE_{T} \, \left(E_{T}^{2} - m_{T}^{2}\right)^{3/2} m_{T} \left( -\frac{F(E_{T} - \sigma\eta_{T})\left[1 - F(E_{T} - \sigma\eta_{T})\right]}{E_{T}} - \frac{F(E_{T} - \sigma\eta_{T})}{E_{T}^{2}} \right).
\end{align}
Substituting the expression for \(\partial \ln \mathcal{Z}/\partial m_{T}\) into \req{eq:magnetization_def} and \req{eq:chain_rule}, we obtain the magnetization
\begin{align}
    \notag
    \mathcal{M} &= -\frac{g|Q|}{2 T}\cdot\frac{2 n_\mathrm{C} V T^{3}}{3 (2\pi)^{2}} \sum_{s}^{\pm1/2} s \sum_{\sigma}^{\pm1} \int_{m_{T}(s)}^{\infty} dE_{T} \, \left(E_{T}^{2} - m_{T}^{2}(s)\right)^{3/2} \\
    \label{eq:mag_form}
    &\times\left[ \frac{F(E_{T} - \sigma\eta_{T})\left(1 - F(E_{T} - \sigma\eta_{T})\right)}{E_{T}} + \frac{F(E_{T} - \sigma\eta_{T})}{E_{T}^{2}} \right].
\end{align}
Much how we expect the free energy to be \(\ln\mathcal{Z}\sim T^{3}\), we see the magnetization is \(\mathcal{M}\sim T^{2}\) via dimensional analysis. This is in agreement to our prior work~\cite{Steinmetz:2023nsc,Steinmetz:2023ucp} where we evaluated the magnetization in the Boltzmann limit. The benefit of expressing the magnetization in the form of \req{eq:mag_form} is that the integrand within the brackets \([\ldots]\) entirely contains the Fermi-Dirac distribution scaled by energy without mass (or magnetic fields) except as a boundary condition on the integration. This makes it suitable for numerical evaluation and comparison to the Boltzmann limit which will be the subject of future work.

%%%%%%%%%%%%%%%%%%%%%%%%
\section{Conclusions}
\label{sec:conclusions}
%%%%%%%%%%%%%%%%%%%%%%%%
In this work we have developed a theoretical framework for evaluating the spin magnetization of a quark-gluon plasma (QGP) under conditions akin to those of the primordial Universe. By employing a grand partition function formalism (see \req{eq:parts}) and rigorously evaluating magnetized Fermi-Dirac integrals, both in their standard and dimensionless forms (Eqs.\,(\ref{eq:dimensionless_partition})--(\ref{eq:dimensionless_partition2})), we derived explicit expressions that capture the dependence of the magnetization on temperature, particle masses, and magnetic field strength. Notably, our analysis shows that the magnetization scales as \(T^2\) (cf. \req{eq:mag_form}), in agreement with the expected thermodynamic behavior of a relativistic gas in three spatial dimensions.

Our estimates indicate that even a modest degree of spin polarization in the light-quark and electron-positron sectors could lead to cosmic magnetic fields of enormous magnitude. Ferromagnetic-like response, if fascilitated by the strong coupling among quarks, would potentially generate fields on the order of \(10^{16}\) Tesla. These values exceed both the critical Schwinger field and the characteristic surface fields of magnetars by several orders of magnitude, underscoring the potential significance of spin magnetization in the primordial plasma.

While the present treatment considers an idealized, free fermion gas, with the magnetization defined in \req{eq:dipole}, it provides a starting point for future investigations incorporate additional physical effects such as QCD interactions, finite volume corrections, and non-equilibrium dynamics. Moreover, exploring the interplay between spin magnetization and other cosmological processes (such as baryogenesis, leptogenesis, and the evolution of large-scale magnetic fields) remains an important avenue for further research.

%%%%%%%%%%%%%%%%%%%%%%%%
\bibliographystyle{sn-aps}
\bibliography{short-note-qcd.bib}
%%%%%%%%%%%%%%%%%%%%%%%%
\end{document}
